\section{Conclusion}
\label{sec-conclusion}
%In this work, we present a novel framework \modelname for generating faithful and reading motive headlines. In addition, we propose a novel dataset \dataset including the title and content of paired fake news and verification news, along with the forward reference type annotation for further research. Our user study found out that verification news suffers from lacking attractiveness, and forward-reference exits in most title with reading motive. Our method applies two additional pre-trained faithfulness scorer and sensation scorer to perform self-verification on both faithfulness and sensationalism during training. Experiment results show that \modelname outperforms baselines in both automatic and human evaluation, which shows its effectiveness. We expect this reading motive headline generator which learns the writing style can help us rewrite boring real news titles and encourage people to read them in order to combat fake news.

We present \modelname, a novel framework for generating faithful
but interesting headlines from a new aspect: forward references. Moreover, we
construct \dataset, a novel dataset that includes the title and content of
pairs of fake and verified news, along with their forward reference types for
further research. Our user study shows that verified news headlines are
relatively boring, and forward references are used in most headlines liked by readers.
Experimental results show that \modelname outperforms all
baselines in both automatic and human evaluations, which demonstrates its
effectiveness in generating attractive but faithful headlines. We expect
\modelname to help rewrite monotonous real news headlines to increase their
exposure rate to help combat fake news.


%We expect this work can help us on rewriting monotonous real news titles, and encourage users' reading motivation on it.

%We expect this reading motive headline generator which learns the writing style can help us rewrite monotonous real news titles and encourage people to read them in order to combat fake news.

% \section*{Limitations}
% Although \modelname shows promising results for re-writing verified news
% headlines to be more attractive while keeping its veracity, there are still
% some limitations: (1)~\modelname is a monolingual model that only
% supports Chinese. It requires three pre-trained scorers. Also, as the FR labels are
% specifically difficult to obtain, it is not easy to implement in
% other languages. (2)~Running the whole framework with a batch size of 16
% takes around 22\,GB GPU memory, mostly because we must load all pre-trained
% models into the GPU. This can be alleviated by using a distilled version as the
% pre-trained model. (3)~On average, \wmodelname generates more faithful
% headlines than other baselines, but it still occasionally produces false
% information or unwanted results. This work is only for academic purposes and is
% not ready for production, and people should double check the results before publishing them.